\documentclass[10pt,a4paper]{amsart}
\usepackage[margin=19mm]{geometry}
\usepackage{amsmath,amssymb,mathtools}
\usepackage[scr=rsfs,cal=euler]{mathalfa}
\usepackage{xfrac}
\usepackage[unicode,hidelinks,bookmarks=false]{hyperref}
\newcommand{\ie}{i.e.\ }
\newcommand{\cf}{cf.\ }
\newcommand{\resp}{respectively\ }
\newcommand{\Subsection}{\subsection}
\providecommand{\nobreakdash}{\nobreak\hbox{-}}
\setlength{\emergencystretch}{2em}
\multlinegap0pt
\usepackage{amssymb}
\newtheorem{proposition}{Proposition}
\usepackage{cleveref}
\crefname{proposition}{Proposition}{Propositions}
\title{Non-reduced components of global nilpotent cones}
\author{David Zhiyuan Bai, David Fang}
\date{Source: arXiv:2604.01331v2 (CC BY 4.0)}
\begin{document}
\maketitle

\noindent\emph{Source and licence.} Non-reduced components of global nilpotent cones. Version arXiv:2604.01331v2 (CC BY 4.0), distributed by arXiv under the Creative Commons Attribution 4.0 International licence, http://creativecommons.org/licenses/by/4.0/, as recorded in the arXiv licence metadata for this exact version and shown on the abstract page https://arxiv.org/abs/2604.01331v2. Full licence notice: \texttt{licenses/arxiv-2604.01331-CC-BY-4.0.txt}.\par\smallskip

\noindent\textit{Editorial scope.} Counted unit: the article's proposition prop:twistedstrata (source0.tex lines 516--518). The article's complete original proof of it is delivered with it (source0.tex lines 549--573, one \texttt{proof} environment of 25 lines, which the article places away from the statement). No mathematics is written, repaired or re-derived: the statement and the proof are the article's own lines, in the article's own notation, and no retained line is substituted or re-flowed.  No further source-local context is needed: every cross-reference of every retained line resolves inside the two delivered spans.  Nothing was omitted from any retained span, so the package carries no omission record.  No retained line contains a citation, so the wrapper carries no bibliography and no bibliography entry.  No retained line contains a figure environment, an image-inclusion command or drawing code, so no image is delivered and the package declares no asset dependency.  Editorial formatting only, in addition: an AMS \texttt{amsart} preamble that restates the article's own declarations and macros used by the retained lines, namely the environments \texttt{proposition} on separate counters, together with the standard packages those lines need.  The retained environments print in this document as Proposition 1 in order of appearance, whereas the article prints the same items with the numbering of its own full document.  The complete native source of the article as distributed in the arXiv e-print for this exact version is bundled unchanged as \texttt{sources/197656/source0.tex}.
\begin{proposition}\label{prop:twistedstrata}
    Suppose $\deg L > 2g-2$. Then $\dim S_m = \dim h_L^{-1}(0)$ if and only if the points of $Z_m$ are generically regular. 
\end{proposition}

\begin{proof}[Proof of \Cref{prop:twistedstrata}]\label{pf:twistedstrata}
    First, note that $\chi(\mathbb{T}_{(E,\phi),w})$ equals
    \[-\sum_{i-j=w}(r_id_j - r_jd_i + r_ir_j(1-g)) + \sum_{i-j=w}(r_id_{j+1} - r_{j+1}d_i + r_ir_{j+1}(l + 1-g)),\]
    where $l=\deg L$.
    In particular, we obtain:
    \begin{align*}
        \sum_{w=-k}^0 \chi(\mathbb{T}_{(E,\phi),w}) &= -\sum_{i}(r_i^2(1-g)) + \sum_{i<j} r_ir_jl\\
        &= \frac{l}2r^2 - \left(\frac{l}2+1-g\right)\sum_ir_i^2.
    \end{align*}

    Since $(E,\phi)$ is stable, we know that
    \[\mathbb{H}^{-1}(C, \mathbb{T}) = \mathbb{H}^0(C, \mathcal H^{-1}(\mathbb{T})) \]
    is one-dimensional and has pure weight $0$.
    On the other hand,
    \[\dim\mathbb H^1(C, \mathbb{T}) = \dim\mathbb H^1(C, \mathcal H^0(\mathbb{T})) =\begin{cases} 1 & L = K_C\\ 0 & \deg L> 2g-2\end{cases}\]
    and $\mathbb H^1(C, \mathcal H^0(\mathbb{T}))$ has pure weight $1$.
    So if $(E,\phi)\in Z_m$, then we can compute
    \[\dim T_{S_m, (E,\phi)} = \sum_{w=-k}^0\dim \mathbb{H}^0(C, \mathbb{T}_{(E,\phi),w}) = \frac{l}2r^2 - \left(\frac{l}2+1-g\right)\sum_ir_i^2 + 1.\]

    When $l >2g-2$, this is bounded above by
    \[\binom{r}{2}l + r(g-1) + 1,\]
    and this bound is obtained when $r_i=1$ for all $i$.
    The very same number is also the relative dimension of $h_L$.
    Therefore $\dim S_m=\dim h_L^{-1}(0)$ precisely when $r_i=1$ for all $i$.
\end{proof}

\end{document}
